\subsection*{6. The Curriculum Is a Dependency Graph, Not a Playlist}

Consider a learner who must diagnose an unreliable AI service. Understanding authorization precedes responsible interpretation of document retrieval. Understanding representative samples precedes credible evaluation. Familiarity with exceptions and logs precedes certain debugging tasks. Some dependencies are logical; others are cognitive, practical, safety-related, or motivational.

Draw a graph of essential capacities, with arrows pointing from prerequisites to later uses. Mark each node as \textbf{must precede}, \textbf{can be developed alongside}, or \textbf{can be supplied when needed}. This prevents the common mistake of treating every interesting topic as a necessary foundation. The fact that mathematics underlies machine learning does not imply that every applied engineering program must begin with measure-theoretic probability. Conversely, avoiding mathematics entirely may leave a program unable to explain evaluation, uncertainty, or optimization. The correct depth comes from the promised performance.

An effective graph is sparse enough to be teachable and honest enough to show hard prerequisites. It allows multiple pathways for different entry populations without pretending that every learner needs the same route. The graph is also an operational asset: when multiple students fail a later task, it identifies possible upstream causes rather than inviting arbitrary extra classes.

\subsection*{7. 3 Useful Shapes of Sequence}

\textbf{Concrete $\to$ abstract $\to$ new concrete} works when learners need initial intuition, formal precision, and later generalization. Use 2 contrasting examples to prevent attachment to incidental features.

\textbf{Worked example $\to$ completion $\to$ independent performance} works for unfamiliar procedures. Guidance should fade; otherwise learners become excellent at following prompts rather than choosing a method.

\textbf{Simple whole $\to$ constituent parts $\to$ more demanding whole} works well for systems and professional judgment. Start by demonstrating the end-to-end job in an accessible form. Study critical components. Reassemble them under altered constraints. The 1st whole should be intelligible, not an exercise in overwhelming novices.

These are not ideological alternatives. A serious course may use all 3 at different moments. The criterion is the structure of the task and the learner's prior knowledge, not the instructor's favorite style.

\subsection*{8. Design a Spiral Rather Than a Parade}

In the parade curriculum, Week 1 teaches topic A, Week 2 topic B, and Week 3 topic C. Once a topic's exam is over, it disappears. The learner becomes familiar with many things and fluent in few. A spiral curriculum deliberately returns to the same governing ideas at increasing depth and under changing conditions.

For a technical studio, \textbf{measurement} might begin as defining a metric, return as selecting an appropriate comparison, return again as detecting leakage, and finally become a decision about whether a system should be released. \textbf{Uncertainty} might begin as probability, return as calibration, return again as error analysis, and become a professional communication obligation. \textbf{Responsibility} might begin as honest citation and later involve handling confidential data or refusing a misleading claim.

The recurring idea must do new work when it returns. Replaying the same slide is not a spiral. The learner needs to retrieve the old structure and use it to solve a more demanding situation. The spacing and retrieval literature supports the value of revisitation, while the specific cadence should be tested against the actual task and retention horizon [S2, S3].

\subsection*{9. Use Mastery Gates Without Worshipping Examinations}

Calendar progress and competence progress can diverge. A critical prerequisite may require an explicit gate: learners proceed when they demonstrate that they can perform an essential task to a defined standard. This can be powerful for safety, mathematical foundations, debugging, and high-stakes practice. But gates impose real costs. Every retry consumes learner and reviewer time, and an overly rigid sequence can trap students on marginal deficiencies unrelated to the final outcome.

Use gates where failure would make later work unintelligible, unsafe, or misleading. For less central material, allow bounded progression with targeted support. Define the retry policy before students struggle. Tell them what evidence is needed, what assistance is permitted, and what happens if the standard remains unmet. \textbf{Rigor is a clear standard coupled with a legitimate path toward it, not a reputation for failing students.}
